% !TeX encoding = UTF-8
% !TeX spellcheck = en_EN-EnglishUnitedKingdom
%% =======================================================
%%		P3C.TeX				v0.3
%%		SCVRIA SOFWARE		Ultimate UOC PEC Repository
%% =======================================================
%%	This file is part of the package P3C.TeX
%%
%%	Author's name: Scvria
%%	File: P3CTeX.code.tex
%%	Version: v0.5 (02/2026)
%%	Description:
%%		Main document class for UOC assignments and projects.
%%		- Owns page layout, typography and global document behaviour.
%%		- Owns student / subject / activity metadata and language selection.
%%		- Delegates feature libraries (UML, property maps, etc.) to pxCORE.
%%
%%	Architectural role:
%%		- P3CTeX.cls never loads standalone libraries (pxUML, pxPRP, ...) directly.
%%			It only talks to the aggregation package `pxCORE` via package options.
%%		- All user-level interfaces that are document-wide (covers, TOC, metadata)
%%			are defined here as LaTeX2e commands.
%%		- expl3 implementation details live in dedicated *.code.tex files or
%%			in library packages; this class only orchestrates them.
%%
%%	This file may be distributed and/or modified:
%%
%%	 1. under the LaTeX Project Public License and/or
%%	 2. under the GNU Public License.
%%
%%	 See the file doc/LICENSE for more details.
%% =======================================================
%%	TODO:	- User level API
%% =======================================================
\ProvidesExplFile{P3CTeX.code.tex}{2026/02/25}{v0.1}
{
	P3CTeX main P3CTeX.cls expl3 code file
}
\bool_new:N \l_pxdoc_language_bool
\bool_set_false:N \l_pxdoc_language_bool
%% =================================
%% ===== P3CTeX KEYS DEFINITION ====
%% =================================
\keys_define:nn {P3CTeX}
{
	cover		.bool_set:N 	=\l_pxdoc_usecover_bool,
		cover		.initial:n 	=	false,
		cover		.default:n 	=	true,
		no-cover	.meta:n 	=	{cover ={ false }},
	toc 		.bool_set:N 	=\l_pxdoc_usetoc_bool,
		toc 		.initial:n 	= 	false,
		toc 		.default:n 	= 	true,
		no-toc		.meta:n 	= 	{toc ={ false }},
	no-plagi	.bool_set:N 	=\l_pxdoc_useplagi_bool,
		no-plagi	.initial:n 	=	false,
		no-plagi	.default:n 	=	true,
	CA		.meta:n	={cover = {false}, toc ={ true }, no-plagi ={true}, en},
	PAC		.meta:n	={cover = {false}, toc ={ true }, no-plagi ={true}, cat},
	PEC		.meta:n	={cover = {false}, toc ={ true }, no-plagi ={true}, es},
	PR		.meta:n	={cover = {true} , toc ={ true }},
	help	.meta:n	={cover = {true} , toc ={ true }, no-plagi ={true}, en, data={help}},
	fake	.meta:n	={cover = {true} , toc ={ true }, no-plagi ={true}, en, data={example}},
	EN		.meta:n	={en},
	ES		.meta:n	={es},
	CAT		.meta:n	={cat},
	cat	.code		={
		\RequirePackage[catalan]{babel}
		\bool_set_true:N \l_pxdoc_language_bool
	},
	es	.code		={
		\RequirePackage[spanish]{babel}
		\bool_set_true:N \l_pxdoc_language_bool
		\_pxcore_set_language_spanish
	},
	en	.code		={
		\RequirePackage[english]{babel}
		\bool_set_true:N \l_pxdoc_language_bool
		\_pxcore_set_language_english
	},
	data 		.code:n	 			= 	{	\keys_set:nn { P3CTeX/data } {#1}	},
	debug 		.bool_set:N 		=		\l_pxcore_debug_bool,
		debug 		.initial:n 		= 		false,
		debug 		.default:n 		= 		true,
	print-doc 	.bool_set:N 		=		\l_pxcore_pxdoc_use_bool,
		print-doc	.initial:n 		= 		false,
		print-doc	.default:n 		= 		true,
	unknown 	.tl_set:N	 		= 	\l__pxdoc_unknown_tl
}
%% =================================
%% == P3CTeX/data KEYS DEFINITION ==
%% =================================
\keys_define:nn { P3CTeX / data}
{
	student			.tl_set:N 		= \l__pxdata_student_tl,
		student			.initial:n	= {Nom~Alumne},
	grade			.tl_set:N		= \l__pxdata_course_tl,
		grade			.initial:n	= {Nom~de~la~Titulació},
	subj-fullname	.tl_set:N		= \l__pxdata_subj_name_tl,
		subj-fullname	.initial:n	= {Nom de la Assignatura},
	subj-shortname	.tl_set:N		= \l__pxdata_subj_short_tl,
		subj-shortname	.initial:n 	= {NmAsg},
	subj-code		.tl_set:N		= \l__pxdata_subj_code_tl,
		subj-code		.initial:n	= {codiAsg},
	ca-fullname		.tl_set:N		= \l__pxdata_activity_tl,
		ca-fullname		.initial:n	= {Nom de l'activitat},
	ca-shortname	.tl_set:N 		= \l__pxdata_act_short_tl,
		ca-shortname	.initial:n	= {NmAct},
	date 			.tl_set:N		= \l__px_date_tl,
		date 			.initial:n	= {\today},
	example	.meta:n	={
		student 		= {Vladimir~Scoriae},
		grade 			= {Software~Engineering~Degree},
		subj-fullname	= {Pràctiques~de~Procrastinació-Avançada~en~\LaTeX},
		subj-shortname	= {PPAL},
		subj-code		= {01312},
		ca-fullname		= {Repte~1:~Introducció~a~\texttt{expl3} },
		ca-shortname	= {PAC1},
		date 			= {1~de~Gener~de~1970}
	},
	help	.meta:n	={
		student 		= {student},
		grade 			= {grade},
		subj-fullname	= {subj-fullname},
		subj-shortname 	= {subj-shortname},
		subj-code 		= {subj-code},
		ca-fullname 	= {ca-fullname},
		ca-shortname 	= {ca-shortname},
		date 			= {date}
	},
	documentation	.meta:n	={
		student 		= {V.Scvria},
		grade 			= {SCVRIA~ SOFWARE},
		subj-fullname	= {Ultimate~ PEC~ Solver~},
		subj-shortname 	= {UOC},
		subj-code 		= {Student~Repository},
		ca-fullname 	= {Universitat~ Oberta~ de~ Catalunya},
		ca-shortname 	= {P3C\TeX},
		date 			= {\today}
	},
	load-data		.code:n	={}
}
\cs_new:Nn \__ifcover:nn {\bool_if:NTF	\l_pxdoc_usecover_bool{#1}{#2}}
\cs_new:Nn \__iftoc:nn	 {\bool_if:NTF	\l_pxdoc_usetoc_bool	{#1}{#2}}
\cs_new:Nn \__ifplagi:nn {\bool_if:NTF	\l_pxdoc_useplagi_bool{#1}{#2}}
%% =================================
%% ======= LOAD P3CTeX DATA ========
%% =================================
\file_input:n {pxTXT.code.tex}%% Load multilanguage text comands
\ProcessKeysOptions{P3CTeX}
\bool_if:NF	\l_pxdoc_language_bool {\RequirePackage[catalan]{babel}} %% Default language setup
\file_input:n {pxGDX.code.tex} %% Load Graphic module
\RequirePackage[\l__pxdoc_unknown_tl]{pxCORE} %% Pass unknown keys to pxCORE
\pxGDX_config:n { %% Pass P3CTeX / data keys to pxCORE
	cover-config={
		text-backtitle	 ={\l__pxdata_course_tl},
		text-uppertitle	 ={\l__pxdata_activity_tl},
		text-title		 ={\l__pxdata_subj_name_tl},
		text-minititle	 ={\l__pxdata_act_short_tl},
		text-subtitle	 ={\l__pxdata_activity_tl},
		text-subsubtitle ={\l__pxdata_subj_code_tl},
		text-author		 ={\l__pxdata_student_tl},
		text-date		 ={\today},
		text-auxtext	 ={ },
	},
	page-config={
		author-name	={\l__pxdata_student_tl},
		doc-name	={\l__pxdata_act_short_tl},
		subj-name	={\l__pxdata_subj_short_tl},
		refcode		={\l__pxdata_subj_code_tl},
		course-name	={\l__pxdata_course_tl}
	}
}
\RenewDocumentCommand{\contentsname}{}{\pxDEFtocname}
\endinput